Our explicit numerical methods—Forward Euler, Improved Euler, and fourth-order Runge-Kutta (RK4)—successfully solved the stiff nonlinear Hodgkin-Huxley equations.
Reliable convergence was achieved with integration steps of $0.07$ to $0.05$ ms, corresponding to a minimum of 14 steps per millisecond.

Simulated recordings were generated by injecting Gaussian noise, amplitude modulation, and physiological artefacts into the ground-truth signal.
A 58th-order Finite Impulse Response (FIR) filter effectively attenuated this noise, with the output root-mean-square (RMS) error scaling linearly with the input noise level.

Spike detection via adaptive thresholding and refractory period enforcement recovered approximately 50\% of true positive spikes.
Detection performance degraded in signal segments containing combined amplitude modulation and ocular artefacts, indicating a limitation of the current threshold-based method.

Several extensions could improve the realism and robustness of the simulation pipeline.
Implementing implicit or exponential integrators would permit larger stable step sizes, reducing computational cost for long simulations.
Realistic, non-stationary noise models such as non-Gaussian interference would better emulate experimental conditions.
Advanced detection algorithms, including template matching or convolutional neural networks, could improve spike recovery in high-artefact regimes.
Finally, validating the pipeline against experimental extracellular recordings would establish its utility for biophysically grounded benchmarking of spike-sorting tools.